\documentclass[9pt,a4paper]{extarticle}
\usepackage[french]{babel}
\usepackage[margin=0.8cm]{geometry}
\usepackage{multicol}
\usepackage{enumitem}
\usepackage{titlesec}
\usepackage{xcolor}
\usepackage{fontspec}

\setmainfont{Segoe UI}

\setlength{\columnsep}{0.4cm}
\setlength{\parindent}{0pt}
\setlength{\parskip}{2pt}

\titleformat{\section}{\normalfont\bfseries\color{blue}}{\thesection}{0.5em}{}
\titleformat{\subsection}{\normalfont\bfseries\small\color{teal}}{\thesubsection}{0.5em}{}
\titlespacing*{\section}{0pt}{5pt}{3pt}
\titlespacing*{\subsection}{0pt}{4pt}{2pt}

\setlist[itemize]{noitemsep, topsep=0pt, leftmargin=*}

\pagestyle{empty}

\begin{document}

\begin{center}
\textbf{\large FICHE DE RÉVISION - SÉCURITÉ DES APPLICATIONS MOBILES}
\end{center}

\begin{multicols}{2}

\section*{ANDROID vs iOS - ARCHITECTURES}

\subsection*{Distribution \& Sandboxing}
\textbf{Android:} Open source, stores alternatifs, sideloading APK possible (risque élevé). SELinux + App Sandbox (UID unique par app). Runtime permissions (Android 6.0+).

\textbf{iOS:} App Store exclusif (revue stricte). App Sandbox isolé + Code Signing obligatoire. Privacy prompts contextuels (iOS 14+).

\subsection*{Fragmentation \& Mises à jour}
\textbf{Android:} 1000+ versions actives, MAJ lentes (fabricant/opérateur).

\textbf{iOS:} >80\% sur iOS 16+ en 6 mois, MAJ OTA rapides Apple direct.

\subsection*{Sécurité Kernel}
\textbf{Android:} Linux Kernel, SELinux enforcing, Verified Boot.

\textbf{iOS:} XNU (Mach+BSD), Mandatory Code Signing, KASLR, Pointer Authentication.

\subsection*{Environnement Sécurisé}
\textbf{Android:} TEE (ARM TrustZone), Android Keystore, StrongBox (Android 9+).

\textbf{iOS:} Secure Enclave (coprocesseur dédié), mémoire isolée, clés crypto.

\section*{SÉCURITÉ OFFENSIVE}

\subsection*{Android - Surface Large}
Sideloading APK, multiplicité OEM, fragmentation patches, root accessible.

\textbf{Attaques:} Malware trojanisé, Banking trojans, Overlay attacks.

\textbf{Coût:} Exploitation moins coûteuse, attaques de masse.

\subsection*{iOS - Surface Restreinte}
Écosystème fermé, signature obligatoire, jailbreak complexe.

\textbf{Attaques:} Exploits kernel zero-day, chaînes Pegasus-like, WebKit.

\textbf{Coût:} Exploits rares, chers, sophistiqués, attaques ciblées.

\section*{MENACES MOBILES}

\textbf{15\%} apps infectées | \textbf{30\%} apps vulnérables réseau | \textbf{+50\%} phishing/an | \textbf{43\%} vol données

\subsection*{Types de Menaces}
\textbf{Malware:} Trojans, Spyware, Ransomware, Banking Trojans.

\textbf{Réseau:} MitM, SSL Stripping, Evil Twin, DNS Spoofing.

\textbf{Phishing:} Smishing (SMS), Fake Apps, Quishing (QR), Vishing.

\textbf{Vol données:} SQLite non chiffré, logs exposés, clipboard, SD card.

\subsection*{Vecteurs d'Attaque}
Apps tierces (>90\% infections), Permissions excessives (90\%+), WiFi publics, Bluetooth/NFC (BlueBorne), Composants exportés.

\section*{ARCHITECTURE ANDROID}

\subsection*{5 Couches}
1. Applications (Gmail, Chrome, apps tierces)

2. Application Framework (Activity/Package Manager, APIs)

3. ART + Native Libraries (OpenSSL, SQLite, WebKit)

4. HAL (Hardware Abstraction Layer)

5. Linux Kernel (drivers, SELinux, namespaces)

\subsection*{Sandbox Android}
\textbf{UID unique:} Range 10000-19999, isolation processus/fichiers.

\textbf{Répertoire privé:} /data/data/ permission 700.

\textbf{IPC sécurisé:} Binder, Intent avec permissions, ContentProvider.

\textbf{SELinux:} MAC enforcing, politiques strictes par domaine.

\subsection*{4 Types de Permissions}
\textbf{Normal:} Auto sans confirmation (INTERNET, BLUETOOTH).

\textbf{Dangerous:} Consentement requis (CAMERA, LOCATION, SMS).

\textbf{Signature:} Même certificat système (INSTALL\_PACKAGES).

\textbf{Special:} Action utilisateur (SYSTEM\_ALERT\_WINDOW).

\subsection*{9 Groupes Permissions Dangereuses}
CALENDAR, CAMERA, CONTACTS, LOCATION, MICROPHONE, PHONE, SENSORS, SMS, STORAGE

\subsection*{Runtime Permissions (Android 6.0+)}
1. checkSelfPermission() - vérifier

2. shouldShowRequestPermissionRationale() - expliquer

3. requestPermissions() - demander

4. onRequestPermissionsResult() - gérer

\textbf{Bonnes pratiques:} Demander au moment de l'utilisation, expliquer l'utilité, gérer le refus gracieusement.

\textbf{Pièges:} Tout demander au lancement, ne pas vérifier (révocable), crasher si refusé.

\subsection*{4 Composants Android}
\textbf{Activity:} Interface utilisateur (exported=false par défaut).

\textbf{Service:} Traitement arrière-plan.

\textbf{BroadcastReceiver:} Écoute événements (BOOT, SMS).

\textbf{ContentProvider:} Partage données inter-apps.

\textbf{Risque exported=true:} Activity lancée par tous, Service exécution arbitraire, ContentProvider accès données, BroadcastReceiver injection intents.

\textbf{Protection:} android:exported="false", android:permission, validation intents, signature-level.

\subsection*{Android KeyStore}
Clés non extractibles, protection TEE/StrongBox.

\textbf{Opérations:} AES-GCM, RSA-OAEP, ECDSA, RSA.

\textbf{Architecture:} App - KeyStore API - daemon - TEE/StrongBox.

\subsection*{TEE \& StrongBox}
\textbf{ARM TrustZone:} Normal World vs Secure World, isolation hardware, SMC.

\textbf{StrongBox (9+):} HSM séparé, CPU/mémoire/stockage isolés, résistance side-channel.

\textbf{Usages:} KeyStore, biométrie, WideVine DRM, Trusted UI.

\subsection*{Stockage Sécurisé}
\textbf{EncryptedSharedPreferences:} AES-256-GCM valeurs, AES-256-SIV clés.

\textbf{EncryptedFile:} AES-256-GCM fichiers.

\textbf{SQLCipher:} AES-256 base SQLite (~10\% perf).

\section*{ARCHITECTURE iOS}

\subsection*{Stack iOS (bas vers haut)}
1. Core OS: Kernel XNU (Mach+BSD)

2. Core Services: Foundation, Core Data, CloudKit

3. Media Layer: AVFoundation, Core Graphics, Metal

4. Cocoa Touch: UIKit, SwiftUI, ARKit

\subsection*{Secure Boot Chain}
1. Boot ROM (iBoot) - code immutable, racine confiance

2. LLB (Low-Level Bootloader) - vérifié signature Apple

3. iBoot (2nd stage) - vérifié par LLB

4. iOS Kernel (XNU) - vérifié par iBoot

5. Kernel Extensions - vérifiées par kernel

\textbf{Échec vérification = arrêt immédiat (recovery/DFU)}

Protection rollback via nonces SHSH.

\subsection*{Secure Enclave}
Coprocesseur ARM dédié, sepOS (microkernel Mach), mémoire chiffrée isolée.

\textbf{Fonctions:} Clés crypto non exportables, Touch ID/Face ID, Apple Pay, TRNG, AES-256.

\subsection*{Code Signing}
\textbf{Certificats:} Development (<100 devices), Ad Hoc (100 max), App Store, Enterprise.

\textbf{Provisioning:} embedded.mobileprovision (cert + App ID + devices + entitlements).

\textbf{Expiration:} 7 jours (gratuit), 1 an (payant).

\textbf{Vérification:} Signature à chaque lancement, intégrité binaire, entitlements.

\subsection*{App Store Review}
1. Analyse auto: malware scan, API privées

2. Review manuelle: guidelines, contenu

3. Tests fonctionnels: crash, comportement

4. Vérification métadonnées

\textbf{Rejets:} API privées, collecte sans consentement, bugs critiques, spam.

\textbf{Limites:} Non infaillible, 24-48h délai, bypass JS hot-patching.

\subsection*{iOS Sandbox \& Permissions}
\textbf{Container:} Bundle/ (lecture seule, signé), Documents/ (iCloud), Library/ (caches), tmp/.

\textbf{Permissions Info.plist:} NSCameraUsageDescription (iOS 10+), NSLocationWhenInUseUsageDescription, NSMicrophoneUsageDescription, NSContactsUsageDescription, NSFaceIDUsageDescription.

\subsection*{Keychain \& Data Protection}
\textbf{Keychain:} AES-256, géré securityd, clés Secure Enclave. Types: GenericPassword, InternetPassword, Certificate, Key, Identity. Accessibilité: kSecAttrAccessibleWhenUnlocked (recommandé).

\textbf{Data Protection - 4 classes:}

\textbf{Complete:} Inaccessible verrouillé (max sécurité).

\textbf{CompleteUnlessOpen:} Ouvert avant verrouillage.

\textbf{CompleteUntilFirstUserAuthentication:} Après 1er unlock (défaut).

\textbf{None:} Toujours accessible (éviter données sensibles).

\section*{POINTS CLÉS}

\subsection*{Android}
Sandbox UID + SELinux | 4 types permissions | Runtime 6.0+ | 9 groupes dangerous | 4 composants | KeyStore TEE/StrongBox | EncryptedSharedPreferences

\subsection*{iOS}
Secure Boot Chain | Code Signing obligatoire | Secure Enclave | App Store Review | Keychain | 4 classes Data Protection | Sandbox container isolé

\subsection*{Comparaison}
\textbf{Android:} Plus ouvert, surface large, attaques masse.

\textbf{iOS:} Plus fermé, surface restreinte, attaques ciblées sophistiquées.

\end{multicols}
\end{document}
